% 実験ガイド v5: エタロン版
\documentclass[12pt,a4paper]{article}

\usepackage{fontspec}
\usepackage{xeCJK}
\setCJKmainfont{Hiragino Mincho ProN}
\setCJKsansfont{Hiragino Kaku Gothic ProN}
\setCJKmonofont{Hiragino Kaku Gothic ProN}
\usepackage[margin=2.5cm]{geometry}
\usepackage{amsmath,amssymb}
\usepackage{hyperref}
\usepackage{tcolorbox}
\usepackage{booktabs}
\usepackage{fancyhdr}

\tcbuselibrary{breakable,skins}

\newtcolorbox{storybox}[1][]{colback=blue!5,colframe=blue!60!black,
  fonttitle=\bfseries\sffamily,title={#1},breakable,arc=3mm}
\newtcolorbox{mathbox}[1][]{colback=green!5,colframe=green!50!black,
  fonttitle=\bfseries\sffamily,title={#1},breakable,arc=3mm}
\newtcolorbox{resultbox}[1][]{colback=yellow!8,colframe=orange!80!black,
  fonttitle=\bfseries\sffamily,title={#1},breakable,arc=3mm}
\newtcolorbox{honestbox}[1][]{colback=gray!8,colframe=gray!50!black,
  fonttitle=\bfseries\sffamily,title={#1},breakable,arc=3mm}
\newtcolorbox{expbox}[1][]{colback=red!5,colframe=red!60!black,
  fonttitle=\bfseries\sffamily,title={#1},breakable,arc=3mm}

\setlength{\parskip}{0.8em}
\setlength{\parindent}{0pt}

\pagestyle{fancy}
\fancyhf{}
\rhead{Wright Brothers}
\lhead{実験ガイド v5}
\cfoot{\thepage}

\begin{document}

\begin{center}
{\Huge\bfseries 素数ミュート実験 v5}\\[0.5cm]
{\Large ファブリ-ペロー・エタロンで偶数モードを消す}\\[0.3cm]
{\Large 室温カシミール実験}\\[1cm]
{\large Wright Brothers Research Program}\\[0.3cm]
{2026年4月}
\end{center}

\vspace{0.5cm}
\tableofcontents
\newpage

%====================================================================
\part{何を測るか}
%====================================================================

\section{一文で}

2枚の板の間の力が、片方の板をエタロンに変えたとき、\textbf{引力から斥力に変わるか}を測る。

\section{もう少し詳しく}

\begin{storybox}[ギターの弦]
ギターの弦には倍音がある: 基本音 ($n=1$)、第2倍音 ($n=2$)、第3倍音 ($n=3$), \ldots

量子力学では、弦を弾かなくても各倍音が微小に「揺れて」いる（ゼロ点揺動）。全倍音の揺れの合計 = 真空エネルギー。

リーマンゼータ関数で計算すると:
\begin{align*}
\text{全モード:} &\quad 1+2+3+4+5+6+\cdots = \zeta(-1) = -\frac{1}{12} \\
\text{奇数のみ:} &\quad 1+3+5+7+9+\cdots = +\frac{1}{12}
\end{align*}

\textbf{偶数倍音を全部消すと、合計の符号が反転する。}
\end{storybox}

\begin{resultbox}[実験の本質]
\begin{center}
\renewcommand{\arraystretch}{1.4}
\begin{tabular}{lcc}
\toprule
配置 & モード & 予測 \\
\midrule
Config A: 両壁が普通のミラー & $1,2,3,4,5,\ldots$ & \textbf{引力} \\
Config B: 片壁がエタロン & $1,3,5,7,9,\ldots$ & \textbf{斥力} \\
\bottomrule
\end{tabular}
\end{center}

AFM（原子間力顕微鏡）で力を測り、符号が変わるかを見る。
\end{resultbox}

%====================================================================
\part{なぜエタロンが偶数を消すか}
%====================================================================

\section{ファブリ-ペロー・エタロンの原理}

\begin{mathbox}[エタロン]
2枚の半透明ミラーの間に薄い空間（厚さ $d$）を挟んだもの。光が中で往復して干渉する。

透過率:
$$T(f) = \frac{1}{1 + F\sin^2(\pi f / \text{FSR})}$$

FSR（自由スペクトル域）$= c/(2d)$。FSR の整数倍の周波数で透過が最大（$T=1$）。それ以外では反射。
\end{mathbox}

\section{エタロンの厚さを a/2 に設定}

キャビティの板間隔を $a$ とする。エタロンの厚さ $d = a/2$。

すると: $\text{FSR} = c/(2 \times a/2) = c/a = 2f_1$

\begin{center}
\renewcommand{\arraystretch}{1.3}
\begin{tabular}{cccc}
\toprule
モード $n$ & 周波数 $f_n = nf_1$ & $\sin^2(n\pi/2)$ & 透過率 \\
\midrule
1 (奇数) & $f_1$ & 1 & 3\%（反射 → \textbf{生存}） \\
2 (偶数) & $2f_1$ & 0 & 100\%（透過 → \textbf{脱出}） \\
3 (奇数) & $3f_1$ & 1 & 3\%（反射 → \textbf{生存}） \\
4 (偶数) & $4f_1$ & 0 & 100\%（透過 → \textbf{脱出}） \\
\vdots & & & \\
\bottomrule
\end{tabular}
\end{center}

\textbf{偶数モードはエタロンを素通りして脱出。奇数モードは反射されてキャビティに閉じ込められる。}

分離は $\sin^2(n\pi/2)$ で\textbf{数学的に完璧}（正確に 0 か 1）。

%====================================================================
\part{予測}
%====================================================================

\section{数値}

\begin{resultbox}[定量的予測]
板間隔 $a = 200$ nm のとき:

\begin{center}
\renewcommand{\arraystretch}{1.3}
\begin{tabular}{lcc}
\toprule
 & Config A（通常ミラー） & Config B（エタロン） \\
\midrule
残るモード & 全て ($1,2,3,\ldots$) & 奇数のみ ($1,3,5,\ldots$) \\
圧力 & $-0.81$ Pa（引力） & $+5.7$ Pa（\textbf{斥力}） \\
\bottomrule
\end{tabular}
\end{center}

力の変化率: \textbf{800\%}。既存のカシミール実験精度（0.1\%）の\textbf{8000倍}大きい。
\end{resultbox}

\begin{honestbox}[これは理想化した計算]
エタロンの反射率を\textbf{完全に}ゼロにはできない。
奇数モードの閉じ込めも完全ではない（透過率 3\%）。
現実のシグナルは理想値より小さくなる。

しかし: シグナルが1\%に劣化しても検出可能（S/N $>$ 80）。
\end{honestbox}

%====================================================================
\part{装置}
%====================================================================

\section{必要なもの}

\begin{expbox}[装置リスト]
\begin{enumerate}
\item \textbf{金コート平面ミラー}（左壁）
    \begin{itemize}
    \item 光学研磨されたシリコン or ガラス基板に金蒸着
    \item 表面粗さ $< 1$ nm RMS
    \item 直径 10 mm 程度
    \end{itemize}

\item \textbf{ファブリ-ペロー・エタロン}（右壁）
    \begin{itemize}
    \item 光学厚さ $d = a/2$（$a = 200$ nm なら $d = 100$ nm）
    \item 半透明ミラー: 誘電体多層膜（$\text{SiO}_2/\text{TiO}_2$）
    \item スペーサー: $\text{SiO}_2$ or $\text{MgF}_2$（厚さ $d$）
    \item 基板: $\text{Si}_3\text{N}_4$ メンブレン（$\sim$50 nm）or 球面レンズ
    \item フィネス $F \geq 10$（偶奇の分離に十分）
    \end{itemize}

\item \textbf{AFM（原子間力顕微鏡）}
    \begin{itemize}
    \item 力の分解能: $< 10$ pN
    \item 距離制御: $< 1$ nm
    \item 大学の共同利用施設で使用可能
    \end{itemize}

\item \textbf{距離制御用ピエゾステージ}
    \begin{itemize}
    \item 板間隔 $a$ を 100 nm 〜 1 $\mu$m でスキャン
    \item 多くの AFM に内蔵
    \end{itemize}
\end{enumerate}
\end{expbox}

\section{エタロンの作り方}

\begin{expbox}[エタロン製作]
\textbf{方法 A: 薄膜蒸着（推奨）}

$\text{Si}_3\text{N}_4$ メンブレン（市販品、50 nm厚）の上に:
\begin{enumerate}
\item 半透明金属膜（Ag, 10 nm）を蒸着 → 半透明ミラー1
\item 誘電体スペーサー（$\text{SiO}_2$, 100 nm）を蒸着
\item 半透明金属膜（Ag, 10 nm）を蒸着 → 半透明ミラー2
\end{enumerate}

これで光学厚さ 100 nm のエタロンが $\text{Si}_3\text{N}_4$ メンブレン上にできる。

蒸着装置は大学のクリーンルームにある。3層を順番に蒸着するだけ。

\textbf{方法 B: 球面エタロン}

球面レンズの表面にエタロン構造をコーティング。
球-平面カシミール実験（Mohideen型）に適合。
AFM の探針にエタロンコートした球を取り付ける。
\end{expbox}

%====================================================================
\part{実験手順}
%====================================================================

\section{Step 1: エタロンの特性確認}

\begin{expbox}[光学テスト]
製作したエタロンの透過スペクトルを分光器で測定。
\begin{itemize}
\item FSR が設計値と一致するか
\item フィネスが 10 以上あるか
\item 偶数モード周波数で透過が最大か
\end{itemize}
これは AFM を使う前に、光学ベンチで確認できる。
\end{expbox}

\section{Step 2: Config A（標準カシミール）}

\begin{expbox}[基準データの取得]
\begin{itemize}
\item 左壁: 金コートミラー
\item 右壁: 金コートミラー（エタロンではない）
\item 板間隔 $a$ を 100 nm 〜 1 $\mu$m でスキャン
\item 各距離での力 $F(a)$ を記録
\item $F \propto -1/a^4$ のカシミール引力カーブを確認
\end{itemize}
これは標準カシミール実験の追試。既知の結果と合うはず。
\end{expbox}

\section{Step 3: Config B（エタロン版 = 素数2ミュート）}

\begin{expbox}[本実験]
\begin{itemize}
\item 左壁: 金コートミラー（同じ）
\item 右壁: エタロン（$d = a/2$ になるよう距離を調整）
\item 板間隔 $a$ をスキャン
\item 各距離での力 $F(a)$ を記録
\end{itemize}

\textbf{注意}: エタロンの厚さ $d$ は固定。$a = 2d$ のとき「素数2ミュート」条件が成立。$a \neq 2d$ では偶奇の分離が不完全になるが、それ自体が有用なデータ。
\end{expbox}

\section{Step 4: 比較}

\begin{expbox}[データ解析]
\begin{enumerate}
\item $a = 2d$ での力を比較:
    \begin{itemize}
    \item Config A: $F_A$（引力、負の値）
    \item Config B: $F_B$（斥力？ 正の値？）
    \end{itemize}
\item $F_B$ の符号を確認:
    \begin{itemize}
    \item $F_B < 0$（引力のまま）→ 標準物理の範囲
    \item $F_B > 0$（斥力に転換）→ \textbf{新発見}
    \end{itemize}
\item $a \neq 2d$ でのデータからエタロンの光学特性を確認
\end{enumerate}
\end{expbox}

%====================================================================
\part{予算}
%====================================================================

\section{費用内訳}

\begin{center}
\renewcommand{\arraystretch}{1.3}
\begin{tabular}{lrl}
\toprule
項目 & 費用 & 備考 \\
\midrule
\multicolumn{3}{l}{\textbf{エタロン製作}} \\
$\text{Si}_3\text{N}_4$ メンブレン & ¥5,000 & 市販品（Norcada, SPI Supplies） \\
蒸着装置使用料（3回） & ¥30,000 & 大学クリーンルーム \\
蒸着材料（Ag, $\text{SiO}_2$） & ¥5,000 & \\
\midrule
\multicolumn{3}{l}{\textbf{ミラー}} \\
金コートシリコン基板 $\times 2$ & ¥10,000 & Thorlabs or 国内業者 \\
\midrule
\multicolumn{3}{l}{\textbf{測定}} \\
AFM 使用料 & ¥50,000 & 大学共同利用（5回） \\
消耗品（カンチレバー等） & ¥10,000 & \\
\midrule
\multicolumn{3}{l}{\textbf{光学テスト}} \\
分光器使用料 & ¥10,000 & \\
\midrule
\multicolumn{2}{r}{\textbf{合計}} & \textbf{¥120,000} \\
\bottomrule
\end{tabular}
\end{center}

\begin{honestbox}[コストの比較]
\begin{center}
\begin{tabular}{lr}
v1-v2 (SQUID, 旧設計): & $\sim$60万円 \\
v3 (トランスモン): & 5〜40万円 \\
v4 (カシミール吸収膜, 撤回): & --- \\
\textbf{v5 (エタロン):} & \textbf{$\sim$12万円} \\
\end{tabular}
\end{center}

v5 は全設計の中で最も安価。しかも室温で動作し、原理的に最も正確。
\end{honestbox}

%====================================================================
\part{技術的課題と対策}
%====================================================================

\section{課題 1: エタロンの品質}

エタロンのフィネスが低いと偶奇の分離が不完全になる。

対策: フィネスは蒸着の膜厚精度で決まる。$\pm 5\%$ の精度があれば $F > 10$ は達成可能。大学のスパッタ装置で十分。

\section{課題 2: 距離 a = 2d の精密制御}

$a = 2d$ から外れると偶奇分離が崩れる。

対策: $a$ をピエゾでスキャンしながら力を記録する。$a = 2d$ の点は力のカーブの変曲点として同定可能。

\section{課題 3: エタロン自身のカシミール力}

エタロンの半透明ミラー間にもカシミール力がかかり、エタロンが変形する可能性。

対策: エタロンの厚さ $d = 100$ nm での内部カシミール圧力は $\sim 10$ Pa。$\text{Si}_3\text{N}_4$ メンブレンの弾性率は十分に高く（$\sim 250$ GPa）、変形は無視できる。

\section{課題 4: 表面力の寄与}

近距離では静電力やファンデルワールス力がカシミール力に重なる。

対策: 標準的なカシミール実験のプロトコル（残留電位の補正、距離校正）を適用。Config A と Config B で同じ補正を行い、差だけを見る。

%====================================================================
\part{共同研究先の探し方}
%====================================================================

\section{必要な設備}

\begin{itemize}
\item AFM（原子間力顕微鏡）: 理学部、工学部にほぼ必ずある
\item 蒸着装置: クリーンルームに標準装備
\item 分光器: 光学実験室に標準装備
\end{itemize}

\textbf{特殊な設備は不要。} 大学の理工学部なら大抵そろっている。

\section{提案の仕方}

\begin{expbox}[提案テンプレート]
「ファブリ-ペロー・エタロンをカシミール実験の片側ミラーとして使い、
モード選択的なカシミール力の変化を測定したい。

エタロンの光学厚さを板間隔の半分に設定すると、
偶数倍音モードはエタロンを透過して脱出し、
奇数倍音モードだけが閉じ込められる。
この条件でカシミール力の符号が反転するかを調べたい。

これはカシミール力のモード選択制御の新しい方法であり、
実験結果は光学コーティング設計や MEMS 応用にも関連する。

エタロンは我々が設計・蒸着する。
AFM の使用を許可いただけないか。」

\textbf{「ゼータ関数」「素数」の話は最初はしない。}
カシミール物理と光学薄膜の研究として提案する。
\end{expbox}

%====================================================================
\part{正直な評価}
%====================================================================

\begin{honestbox}[この実験でわかること]

\textbf{もし斥力が観測されたら}:
\begin{itemize}
\item 偶数モードの除去が真空エネルギーの符号を反転させた最初の証拠
\item ゼータ関数の算術構造が物理的真空に入っている示唆
\item カシミール力のモード選択制御の実証
\item Nature / Science 級の発見
\end{itemize}

\textbf{もし引力のままだったら}:
\begin{itemize}
\item エタロンの不完全さ（有限フィネス、反射ロス）が原因の可能性
\item 標準 Lifshitz 理論で説明可能なら、新物理なし
\item ただし: 「エタロンによるモード選択カシミール力」自体は新しいデータ
\item 光学薄膜とカシミール効果の分野で論文になりうる
\end{itemize}

\textbf{偶然の確率}: この実験は cos $\theta_W$ = reg(Q($\sqrt{2}$)) とは独立。数学的定理（符号反転定理）に基づく。「数値の一致」ではなく「力の符号」を見るので、偶然の問題は小さい。
\end{honestbox}

\vfill

\begin{center}
\rule{10cm}{0.4pt}\\[0.5cm]
{\Large\bfseries エタロンを通り抜けるか、反射されるか。}\\[0.2cm]
{\Large\bfseries 偶数モードと奇数モードの運命を分ける鏡。}\\[0.3cm]
{\small Wright Brothers Research Program, 2026}
\end{center}

\end{document}
